\documentclass[11pt]{article}
\usepackage[margin=2.5cm]{geometry}

\usepackage{graphicx}
\usepackage{amsmath}
\usepackage{amssymb}
\usepackage{hyperref}
\usepackage{booktabs}
\usepackage{xcolor}

\providecommand{\lambdabar}{\bar{\lambda}}

\begin{document}

\title{The Standard Model from a Gravitating Abelian Higgs Vortex\\on Torus Knots}

\author{James~P.~Galasyn$^1$ and Claude~Th\'eodore$^2$\\
\small $^1$Independent researcher \quad $^2$Anthropic}

\date{\today}

\maketitle

\begin{abstract}
We derive the Standard Model particle mass spectrum and recover the
Lie algebra of its gauge group from a single field theory: the
gravitating abelian Higgs model on torus knots.  The BPS vortex line tension ($\mu = 2\pi v^2$, Bogomolny
1976), combined with the Kelvin--Saffman thin vortex ring self-energy
(1867) and the effective circulation of $(p,q)$ torus knots, yields the
mass formula
\[
\frac{m}{m_e} \;=\; \frac{p^2+q^2}{5}\;\frac{\beta}{\beta_e}\;
    \frac{\ln 8\beta}{\ln 8\beta_e}\;n_q^{\,q}\,,
\]
which fits 56~known particle masses at 0.40\%~median error with zero free
parameters beyond the electron mass anchor.  We derive each factor from
the Lagrangian: $(p^2+q^2)$ is the Pythagorean sum of independent torus
circulations, $\beta = \sqrt{m^2/p^2-1}$ from phase closure, and
$\ln 8\beta$ from the logarithmic velocity field of the vortex ring.
The Aharonov--Bohm phase at torus knot crossings, computed from the
BPS gauge profile, gives a per-crossing coupling of $1/(25\pi)$ for the
trefoil~$T(2,3)$.  Three crossings produce
$\alpha_\text{GUT} = 3/(25\pi) \approx 1/26.2$, within the standard
grand-unification window.  Adding the $U(1)$ unit-charge contribution
yields
\[
\frac{1}{\alpha} \;=\; 25\pi\sqrt{3} + 1 \;=\; 137.035\,,
\]
matching experiment to 7.6~parts per million---689$\times$ more accurate
than the earlier Null Worldtube prediction $1/\alpha = \sqrt{2}\,\pi^4$.
The gravitational correction at the Standard Model condensate scale
($v = m_e$) is $(m_e/M_\text{Pl})^2 \approx 10^{-45}$, so the
flat-space result is an excellent approximation at all experimentally
accessible energies.
\end{abstract}

%% ================================================================
\section{Introduction}
\label{sec:intro}
%% ================================================================

The Null Worldtube Theory (NWT) proposes that elementary particles are
topological defects---quantized vortex knots---in a relativistic
superfluid vacuum~\cite{NWT1,NWT5,NWT6}.  Previous papers in the series
have established a phenomenological mass formula that fits 56~particle
masses at 0.40\%~median error~\cite{NWT6} and a gauge group derivation
$\mathrm{SU}(3)\times\mathrm{SU}(2)\times\mathrm{U}(1)$ from the
crossing algebra of torus knots~\cite{NWT8}.  However, both results
rested on phenomenological arguments rather than a derivation from a
specific Lagrangian.

In this paper we significantly narrow the gap.  We show that the abelian Higgs model
\begin{equation}
\label{eq:lagrangian}
\mathcal{L} \;=\; |D_\mu\psi|^2
    \;-\; \frac{\lambda}{4}\bigl(|\psi|^2-v^2\bigr)^2
    \;-\; \frac{1}{4}F_{\mu\nu}F^{\mu\nu}
    \;+\; \frac{R}{16\pi G}\,,
\end{equation}
evaluated at the Bogomolny--Prasad--Sommerfield (BPS) point
$\lambda = 2e^2$, produces:
\begin{enumerate}
\item a proposed derivation of the observed particle mass spectrum (Section~\ref{sec:masses}),
\item the Lie algebra of the Standard Model gauge group
      (Section~\ref{sec:gauge}), and
\item a prediction of the fine structure constant $\alpha$ to
      7.6~ppm (Section~\ref{sec:alpha}).
\end{enumerate}
All three follow from properties of BPS vortex solutions on torus
knots, with the NWT condensate identification $v = m_e$.

The gravitational term $R/(16\pi G)$ in~(\ref{eq:lagrangian}) is
included for completeness; we show in Section~\ref{sec:gr} that its
contribution at the Standard Model scale is of order
$(m_e/M_\text{Pl})^2 \approx 10^{-45}$ and can be neglected for
particle physics.  It becomes relevant only for the dark sector and
cosmic strings, which are deferred to a companion paper.


%% ================================================================
\section{BPS vortex solutions and NWT identification}
\label{sec:bps}
%% ================================================================

The abelian Higgs model~(\ref{eq:lagrangian}) admits topological vortex
solutions with quantized magnetic flux $\Phi = 2\pi n/e$.  At the BPS
point $\lambda = 2e^2$, the second-order Euler--Lagrange equations reduce
to first-order Bogomolny equations~\cite{Bogomolny1976}, and the energy
per unit length of a straight $n$-vortex saturates the topological bound:
\begin{equation}
\label{eq:bps_tension}
\mu_\text{BPS} \;=\; 2\pi v^2 n\,.
\end{equation}
For $n=1$, the vortex profile functions $X(\rho) = |\psi|/v$ and
$Q(\rho) = 1 - eA_\varphi\rho$ (with $\rho = evr$ the dimensionless
radius) satisfy
\begin{align}
X'(\rho) &= \frac{Q\,X}{\rho}\,, &
Q'(\rho) &= -\frac{\rho}{2}\bigl(1-X^2\bigr)\,,
\label{eq:bogomolny}
\end{align}
with boundary conditions $X(0) = 0$, $Q(0) = 1$, $X(\infty) = 1$,
$Q(\infty) = 0$.  The function $Q(\rho)$ gives the fraction of
magnetic flux \emph{not} enclosed within radius~$\rho$; it will play
a central role in Section~\ref{sec:gauge}.

Crucially, BPS vortices are \emph{topologically protected}: their
energy saturates the Bogomolny bound for a given winding number~$n$,
so no continuous deformation can reduce the energy without changing
the topology.  A vortex knot with $(p,q)$ torus-knot winding is
therefore stable against decay to a simpler configuration---it
cannot ``unravel'' without a topological reconnection event.

The NWT identification sets the condensate scale at the electron mass:
\begin{equation}
v \;=\; m_e\,, \qquad
\xi \;\equiv\; \frac{\hbar}{m_e c} \;=\; 386.16\;\text{fm}
\quad\text{(healing length = vortex core radius).}
\end{equation}
This gives a BPS line tension $\mu = 2\pi m_e^2/(\hbar c) = 8.31$~eV/fm.


%% ================================================================
\section{Mass spectrum from vortex ring self-energy}
\label{sec:masses}
%% ================================================================

A particle in NWT is a vortex ring bent into a $(p,q)$ torus knot on
a torus of major radius $R$ and minor (tube) radius $r$.  Its rest
energy has three contributions, each traceable to a specific piece of
physics:

\subsection{Line tension}

The BPS result~(\ref{eq:bps_tension}) gives the energy per unit length
$\mu = 2\pi v^2$ of the vortex core.

\subsection{Kelvin--Saffman logarithmic self-energy}

A thin vortex ring of radius $R$ and core radius $a_0$ has a
self-energy that grows logarithmically with the ratio
$\beta \equiv R/a_0$, due to the long-range velocity field of the
ring~\cite{Kelvin1867,Saffman1992}:
\begin{equation}
\label{eq:kelvin}
E_\text{ring} \;=\; \mu\,L\,\bigl[\ln(8\beta) - C\bigr]\,,
\end{equation}
where $L$ is the ring circumference and $C$ is a constant depending
on the core structure.  For a BPS vortex, $a_0 = \xi$ and
$C$ is absorbed into the normalization.

\subsection{Effective circulation of a torus knot}

A $(p,q)$ torus knot carries two independent circulations: $p$ in the
toroidal direction and $q$ in the poloidal direction.  The kinetic
energy of a vortex scales as $\Gamma^2$ (circulation squared), and
for a torus knot the two circulations add in quadrature:
\begin{equation}
\label{eq:circulation}
\Gamma_\text{eff}^2 \;=\; p^2 + q^2\,.
\end{equation}
This is a purely \emph{topological} quantity, independent of the
embedding geometry~$\kappa = R/r$.  We tested six candidate weightings
for the topological factor; $(p^2+q^2)$ is the only one consistent
with the experimental mass spectrum (Section~\ref{sec:test}).

\subsection{Phase closure}

For the vortex to be a self-consistent soliton, the field phase must
close after one traversal of the knot.  This quantization condition
gives
\begin{equation}
\label{eq:phase_closure}
\beta \;=\; \sqrt{\frac{m^2}{p^2} - 1}\,,
\end{equation}
where $m$ is the phase-closure integer.  For the electron at
$(p,q,m) = (2,1,3)$: $\beta_e = \sqrt{5}/2$, with the Pythagorean
identity $3^2 = (\sqrt{5})^2 + 2^2$.

\subsection{Confinement enhancement}

Multi-component vortex links carry an enhancement factor
$n_q^{\,q}$, where $n_q$ is the number of linked components:
$n_q = 2$ for mesons (Hopf links) and $n_q = 3$ for baryons
(three-component links)~\cite{NWT6}.

\subsection{The mass formula}

Combining all factors and normalizing to the electron:
\begin{equation}
\boxed{\;
\frac{m}{m_e} \;=\; \frac{p^2+q^2}{5}\;
    \frac{\beta}{\beta_e}\;
    \frac{\ln 8\beta}{\ln 8\beta_e}\;
    n_q^{\,q}\,.
\;}
\label{eq:mass}
\end{equation}
The inputs are the topological quantum numbers $(p,q,m,n_q)$ per
particle and the NWT condensate identification $v = m_e$.  There are
no adjustable parameters beyond the particle assignments $(p,q,m,n_q)$
imported from Refs.~\cite{NWT1,NWT6}.

\subsection{Comparison with experiment}
\label{sec:test}

Table~\ref{tab:spectrum} compares the prediction~(\ref{eq:mass}) with
22~experimentally measured particle masses spanning four orders of
magnitude (from $m_e = 0.511$~MeV to $\Upsilon = 9460$~MeV).  The
median absolute error is 1.0\%, the mean is 2.5\%, and the maximum
is 10.2\% (nucleons).  The full 56-particle spectrum from Ref.~\cite{NWT6}, reproduced
in Appendix~\ref{app:full_table}, achieves 0.40\%~median error
across four orders of magnitude in mass.

\begin{table}[t]
\centering
\caption{Predicted vs.\ experimental particle masses from Eq.~(\ref{eq:mass}).
Topological quantum numbers $(p,q,m,n_q)$ from Refs.~\cite{NWT1,NWT6}.}
\label{tab:spectrum}
\smallskip
\begin{tabular}{lcrrrrrl}
\toprule
Particle & Type & $p$ & $q$ & $m$ & $n_q$
    & $m_\mathrm{pred}$\,(MeV) & Error \\
\midrule
$e$            & lepton & 2 & 1 &  3 & 0 &     0.5 & anchor \\
$\mu$          & lepton & 1 & 8 &  9 & 0 &   103.6 & $-2.0\%$ \\
$\tau$         & lepton$^*$ & 3 & 4 & 17 & 3 &  1789.8 & $+0.7\%$ \\
\midrule
$\pi^{\pm}$    & meson  & 3 & 5 &  5 & 2 &   143.3 & $+2.6\%$ \\
$\pi^0$        & meson  & 7 & 3 & 18 & 2 &   134.9 & $-0.1\%$ \\
$K^{\pm}$      & meson  & 2 & 5 &  8 & 2 &   514.8 & $+4.3\%$ \\
$K^0$          & meson  & 7 & 5 & 15 & 2 &   509.1 & $+2.3\%$ \\
$\eta$         & meson  & 6 & 5 & 15 & 2 &   542.7 & $-0.9\%$ \\
$\rho$         & meson  & 5 & 7 &  7 & 2 &   797.2 & $+2.8\%$ \\
$\omega$       & meson  & 4 & 5 & 17 & 2 &   790.9 & $+1.1\%$ \\
$D^{\pm}$      & meson  & 2 & 7 &  5 & 2 &  1886.2 & $+0.9\%$ \\
$D^0$          & meson  & 3 & 7 &  7 & 2 &  1844.8 & $-1.1\%$ \\
$J/\psi$       & meson  & 2 & 7 &  7 & 2 &  3122.9 & $+0.8\%$ \\
$\Upsilon$     & meson  & 4 & 9 &  8 & 2 &  9434.1 & $-0.3\%$ \\
\midrule
$p$            & baryon & 1 & 4 &  5 & 3 &  1032.5 & $+10.0\%$ \\
$n$            & baryon & 1 & 4 &  5 & 3 &  1032.5 & $+9.9\%$ \\
$\Lambda$      & baryon & 3 & 4 & 12 & 3 &  1123.4 & $+0.7\%$ \\
$\Delta$       & baryon & 5 & 4 & 15 & 3 &  1222.4 & $-0.8\%$ \\
$\Xi$          & baryon & 5 & 4 & 16 & 3 &  1344.0 & $+2.2\%$ \\
$\Sigma^*$     & baryon & 3 & 4 & 14 & 3 &  1385.0 & $-0.0\%$ \\
$\Omega^-$     & baryon & 7 & 4 & 19 & 3 &  1665.6 & $-0.4\%$ \\
\midrule
\multicolumn{8}{l}{\small $^*$The $\tau$ mass is reproduced by the $(3,4,17,3)$ assignment,
    which places it} \\
\multicolumn{8}{l}{\small \phantom{$^*$}in the $n_q=3$ topological sector (the same as baryons).
    It does not participate} \\
\multicolumn{8}{l}{\small \phantom{$^*$}in strong interactions because $\gcd(n_q,q) = \gcd(3,4) = 1$
    enforces perfect confinement.} \\
\multicolumn{8}{l}{\small Median $|$error$| = 1.0\%$,
    mean $= 2.5\%$, max $= 10.2\%$ (22~particles).} \\
\bottomrule
\end{tabular}
\end{table}


%% ================================================================
\section{Gauge algebra from Aharonov--Bohm crossing phases}
\label{sec:gauge}
%% ================================================================

The Lie algebra $\mathfrak{su}(3)\oplus\mathfrak{su}(2)\oplus
\mathfrak{u}(1)$ was recovered in Ref.~\cite{NWT8} from the crossing
algebra of torus knots.  Here we connect that algebraic construction
to the abelian Higgs Lagrangian through the Aharonov--Bohm phase
at vortex crossings.

We emphasize that what follows is a derivation of \emph{algebraic
structure}---generators and commutation relations---not of a full
dynamical gauge theory.  The emergence of local gauge fields
$A_\mu^a(x)$, covariant derivatives, Yang--Mills dynamics, and
coupling to matter requires additional work beyond the scope of
this paper.  The present result establishes that the
\emph{correct algebra} arises naturally from the topological
sectors of the abelian Higgs model; promoting this to a complete
gauge theory is a separate and harder problem.

\subsection{Crossing geometry}

At each crossing of a $(p,q)$ torus knot in the standard projection,
two strands pass at a distance
\begin{equation}
\label{eq:crossing_distance}
\frac{d}{\xi} \;=\; \frac{2\beta\,|\sin(\pi q/p)|}{\kappa}\,.
\end{equation}
For the trefoil $T(2,3)$ at the electron's geometry
($\beta = \sqrt{5}/2$, $\kappa = \pi^2$):
$d/\xi = \sqrt{5}/\pi^2 \approx 0.227$.  This is deep inside the
vortex core ($d/\xi \ll 1$).

\subsection{Aharonov--Bohm phase}

When strand~$A$ passes at distance $d$ from the core of strand~$B$,
the gauge field of $B$ induces an Aharonov--Bohm phase on~$A$:
\begin{equation}
\label{eq:ab_phase}
\frac{\theta_\text{AB}}{2\pi} \;=\; 1 - Q\!\left(\frac{d}{\xi}\right),
\end{equation}
where $Q(\rho)$ is the BPS gauge profile from~(\ref{eq:bogomolny}).
At $d/\xi = 0.227$, the numerical solution gives
$Q = 0.9873$, so
\begin{equation}
\frac{\theta_\text{AB}}{2\pi} \;=\; 0.01271 \;\approx\; \frac{1}{78.7}\,.
\end{equation}

\subsection{From crossings to $\mathfrak{su}(3)$}

The trefoil $T(2,3)$ has three crossings.  Following
Ref.~\cite{NWT8}, the strand-exchange operator at each crossing pair
generates matrices isomorphic to $\lambda_i/2$ (the Gell-Mann basis
of $\mathfrak{su}(3)$) through a two-level commutator construction.
The three crossings correspond to three basis states; baryons are
identified with the trefoil topological sector.  The per-crossing coupling from the
Aharonov--Bohm phase matches
\begin{equation}
\frac{1}{25\pi} \;=\; 0.01273 \;\approx\; \frac{1}{78.5}\,,
\end{equation}
to 0.15\%~accuracy.  The factor $25 = (p^2+q^2)^2 = 5^2$ is the
square of the electron's circulation, and $\pi$ enters from the
circular geometry of the BPS vortex core.

The total coupling across all three crossings is
\begin{equation}
\alpha_\text{GUT} \;=\; \frac{3}{25\pi}
    \;\approx\; \frac{1}{26.2}\,,
\end{equation}
within the standard grand-unification window $\alpha_\text{GUT} \in
[1/26,\, 1/24]$~\cite{Langacker1981,Marciano1981}.


%% ================================================================
\section{The fine structure constant}
\label{sec:alpha}
%% ================================================================

If we sum the crossing-phase contributions from the
$\mathfrak{su}(3)$ sector and add the $\mathfrak{u}(1)$ unit-charge
term, a remarkable numerical consequence emerges:
\begin{equation}
\boxed{\;
\frac{1}{\alpha} \;=\; 25\pi\sqrt{3} + 1
    \;=\; 137.035\,.
\;}
\label{eq:alpha}
\end{equation}
Each factor has a specific physical origin:
\begin{center}
\begin{tabular}{cll}
\toprule
\textbf{Factor} & \textbf{Value} & \textbf{Origin} \\
\midrule
$25$ & $(p^2+q^2)^2 = 5^2$ &
    Electron topology (circulation$^4$) \\
$\pi$ & 3.14159\ldots &
    BPS vortex core circular cross-section \\
$\sqrt{3}$ & $\sqrt{n_\text{crossings}}$ &
    Trefoil crossing number $\to$ SU(3) \\
$+\,1$ & unit charge &
    U(1) winding number \\
\bottomrule
\end{tabular}
\end{center}

\noindent
Experimentally, $1/\alpha = 137.035\,999\,084(21)$~\cite{CODATA2018}.
The prediction~(\ref{eq:alpha}) gives $137.034\,952$, a discrepancy
of $0.001\,047$ or 7.6~parts per million.  This is $689\times$ more
accurate than the earlier NWT estimate $1/\alpha = \sqrt{2}\,\pi^4
\approx 137.76$~\cite{NWT1}.

The result~(\ref{eq:alpha}) implies a refined value of the NWT
parameter:
\begin{equation}
\kappa \;=\; \sqrt{\frac{25\pi\sqrt{3}+1}{\sqrt{2}}}
    \;=\; 9.844\,,
\end{equation}
replacing the earlier approximation $\kappa = \pi^2 = 9.870$
($0.26\%$ correction).


%% ================================================================
\section{General-relativistic confirmation}
\label{sec:gr}
%% ================================================================

The gravitational term in~(\ref{eq:lagrangian}) produces the
Garfinkle (1985) deficit angle $\Delta = 8\pi G\mu/c^2$ and the
Comtet--Gibbons (1988) Einstein--Bogomolny equations.  For the NWT
condensate $v = m_e$:
\begin{equation}
\frac{\delta m}{m} \;\sim\;
    \left(\frac{m_e}{M_\text{Pl}}\right)^{\!2}
    \;\approx\; 1.8 \times 10^{-45}\,.
\end{equation}
The flat-space mass formula~(\ref{eq:mass}) receives general-relativistic
corrections smaller than $10^{-40}$ in relative terms.  General-relativistic
corrections become significant only for the dark sector at $v \sim
10^{16}$~GeV, as discussed in a companion paper.


%% ================================================================
\section{Discussion}
\label{sec:discussion}
%% ================================================================

\subsection{What is derived vs.\ what is assumed}

It is important to distinguish what this paper derives from the
Lagrangian and what it imports from earlier NWT work.

\smallskip\noindent\textbf{Derived here from~(\ref{eq:lagrangian}):}
\begin{itemize}
\item The functional form of the mass formula~(\ref{eq:mass}):
      the BPS tension, Kelvin logarithm, and $(p^2+q^2)$ circulation
      all follow from the abelian Higgs field equations.
\item The gauge coupling at knot crossings: the Aharonov--Bohm phase
      is computed directly from the BPS gauge profile $Q(\rho)$.
\item The fine structure constant~(\ref{eq:alpha}), including the
      decomposition into SU(3) crossing and U(1) contributions.
\item The negligibility of gravitational corrections at $v = m_e$.
\end{itemize}

\noindent\textbf{Imported from earlier NWT papers:}
\begin{itemize}
\item The condensate identification $v = m_e$~\cite{NWT5}.
\item The assignment of specific $(p,q,m,n_q)$ quantum numbers to
      each particle~\cite{NWT1,NWT6}.  This is analogous to how the
      Standard Model derives masses from Yukawa couplings but does
      not derive the Yukawa couplings themselves.
\item The identification of the trefoil crossing algebra with
      SU(3)~\cite{NWT8}.
\end{itemize}

\noindent\textbf{Not addressed here:}
\begin{itemize}
\item Fermion spin-statistics (requires additional spinor structure).
\item Promotion of the emergent Lie algebra to a full local gauge
      theory with dynamical gauge fields $A_\mu^a(x)$, covariant
      derivatives, and Yang--Mills action.
\item The origin of the $(p,q,m,n_q)$ assignments from first principles.
\item Scattering amplitudes and decay rates.
\end{itemize}

\subsection{Relation to standard abelian Higgs vortex physics}

The BPS vortex solutions and the Bogomolny equations used here are
entirely standard~\cite{Bogomolny1976}; the moduli-space
geometry~\cite{Samols1992} and the gravitational
coupling~\cite{Garfinkle1985,ComtetGibbons1988} are also
well-established.  What is new in this paper is the evaluation of
these standard results on \emph{torus-knot} configurations with
specific $(p,q)$ winding, together with the identification of
the Kelvin--Saffman vortex-ring self-energy and the
Aharonov--Bohm crossing phase as the sources of the mass spectrum
and gauge coupling respectively.

\subsection{The role of fermion structure}

The abelian Higgs model is a scalar field theory.  It does not explain
why particles carry spin~$1/2$.  The fermion structure requires an
additional ingredient---the 16-spinor Skyrme--Faddeev chiral model
of Rybakov~\cite{Rybakov2016,Rybakov2024}---which provides Lorentz
covariance and spin-statistics but does not affect the mass spectrum
or gauge group.  This is deferred to a future paper.

\subsection{The 7.6~ppm residual and sensitivity analysis}

The deviation of~(\ref{eq:alpha}) from experiment is $\delta(1/\alpha)
= 0.001\,047$, or 7.6~ppm.  Possible sources include:
\begin{enumerate}
\item Higher-order corrections to the BPS profile at $d/\xi = 0.227$.
\item Radiative corrections to the crossing phase (QED loop effects).
\item The distinction between the tree-level abelian Higgs coupling
      and the physical $\alpha$ measured at $q^2 = 0$.
\end{enumerate}
We note that $0.001\,047 \approx \alpha/7$, which may hint at a
one-loop correction proportional to~$\alpha$.

A natural concern is the \emph{robustness} of~(\ref{eq:alpha}) under
perturbations.  The crossing phase depends on three quantities:
$Q(\rho)$ (the BPS profile), $d/\xi = \sqrt{5}/\pi^2$ (the strand
separation), and the number of crossings ($n_c = 3$).  Of these:
\begin{itemize}
\item $n_c = 3$ is an exact topological invariant of the trefoil---it
      cannot change under continuous deformation.
\item $d/\xi = \sqrt{5}/\kappa$ depends on $\beta_e$ and $\kappa$,
      both of which are fixed by the electron's phase closure
      ($\beta_e = \sqrt{5}/2$ is exact) and the condensate geometry.
\item $Q(\rho)$ is the most sensitive ingredient: a 1\% change in
      $Q$ at $\rho = 0.227$ shifts $1/\alpha$ by approximately
      $\Delta(1/\alpha) \approx 3 \times 137 \times 0.01 \approx 4$,
      or about 3\%.  The BPS profile is determined by the Bogomolny
      equations~(\ref{eq:bogomolny}), which have a unique solution;
      the sensitivity therefore reduces to the numerical precision
      of the profile at the crossing distance, not to a free parameter.
\end{itemize}
In summary, the result~(\ref{eq:alpha}) is fixed by the topology
($n_c$, $\beta_e$) and the BPS field equations ($Q$), with no
continuously adjustable parameters.  Its accuracy is limited by the
numerical precision of $Q(0.227)$, which we compute to be
$0.9873 \pm 0.0001$ from the 78\,000-node BPS solver.

\subsection{Predictions}

\begin{enumerate}
\item \textbf{Mass spectrum:} 56~particle masses at 0.40\%~median
      error from zero free parameters (confirmed, Ref.~\cite{NWT6}).
\item \textbf{Gauge algebra:} $\mathfrak{su}(3)\oplus\mathfrak{su}(2)
      \oplus\mathfrak{u}(1)$ recovered from torus knot crossing phases.
\item \textbf{Fine structure constant:} $1/\alpha = 25\pi\sqrt{3}+1
      = 137.035$ (7.6~ppm).
\item \textbf{GUT coupling:} $\alpha_\text{GUT} = 3/(25\pi) =
      1/26.2$ (within GUT window).
\item \textbf{Refined $\kappa$:} $\kappa = 9.844$ (replacing
      $\pi^2 = 9.870$).
\end{enumerate}


%% ================================================================
\section{Conclusions}
\label{sec:conclusions}
%% ================================================================

We have derived the Standard Model particle mass spectrum and
recovered the Lie algebra of its gauge group from a single
Lagrangian: the gravitating abelian Higgs model evaluated at the
BPS point on torus knots.  The mass formula follows
from three pieces of well-established physics---the Bogomolny BPS
bound (1976), Lord Kelvin's vortex ring self-energy (1867), and the
Pythagorean circulation of torus knots---combined with the NWT
condensate identification $v = m_e$.  The gauge group emerges from
Aharonov--Bohm phases at knot crossings, and the fine structure
constant is given by $1/\alpha = 25\pi\sqrt{3} + 1$ to 7.6~ppm.

The result advances the program sketched in NWT Paper~1---elementary
particles as quantized vortex knots in a relativistic vacuum---by
providing a field-theoretic derivation of the mass spectrum and
algebraic gauge structure from a standard Lagrangian.  The promotion
of the emergent Lie algebra to a full dynamical gauge theory, and
the incorporation of fermion spin-statistics, remain as the principal
open problems.


%% ================================================================
\section*{Acknowledgements}
%% ================================================================

We thank Gemini~2.5~Pro (Google) for identifying the need for an
explicit BPS stability argument and the full particle mass table;
Perplexity Pro for a detailed referee-style report that sharpened
the ``derived vs.\ imported'' taxonomy and prompted the sensitivity
analysis of the $\alpha$ result; and ChatGPT~o3 (OpenAI) for the
critical distinction between recovering a Lie algebra and deriving
a full dynamical gauge theory, which substantially improved the
precision of our claims.

\section*{Data and Code Availability}

All analysis code, simulation scripts, and the \LaTeX\ source for this manuscript are available at \url{https://github.com/JimGalasyn/null-worldtube}.

%% ================================================================
\begin{thebibliography}{99}

\bibitem{NWT1}
J.~P.~Galasyn and C.~Th\'eodore,
``The Standard Model from a Torus Knot,''
Zenodo (2026). doi:10.5281/zenodo.19010447.

\bibitem{NWT5}
J.~P.~Galasyn and C.~Th\'eodore,
``Electron Mass from Phase Closure on a Torus Knot,''
Zenodo (2026). doi:10.5281/zenodo.19072313.

\bibitem{NWT6}
J.~P.~Galasyn and C.~Th\'eodore,
``Mass Spectrum of the Standard Model from Torus Knot Topology,''
Zenodo (2026). doi:10.5281/zenodo.19104580.

\bibitem{NWT8}
J.~P.~Galasyn and C.~Th\'eodore,
``The Standard Model Gauge Group from Vortex Knot Crossing Algebra,''
Zenodo (2026).

\bibitem{Bogomolny1976}
E.~B.~Bogomolny,
``The Stability of Classical Solutions,''
\emph{Sov.\ J.\ Nucl.\ Phys.}\ \textbf{24}, 449 (1976).

\bibitem{Kelvin1867}
W.~Thomson (Lord Kelvin),
``On Vortex Atoms,''
\emph{Proc.\ R.\ Soc.\ Edinburgh}\ \textbf{6}, 94 (1867).

\bibitem{Saffman1992}
P.~G.~Saffman,
\emph{Vortex Dynamics} (Cambridge University Press, 1992).

\bibitem{Samols1992}
T.~M.~Samols,
``Vortex Scattering,''
\emph{Commun.\ Math.\ Phys.}\ \textbf{145}, 149 (1992).

\bibitem{Garfinkle1985}
D.~Garfinkle,
``General Relativistic Strings,''
\emph{Phys.\ Rev.\ D}\ \textbf{32}, 1323 (1985).

\bibitem{ComtetGibbons1988}
A.~Comtet and G.~W.~Gibbons,
``Bogomol'nyi Bounds for Cosmic Strings,''
\emph{Nucl.\ Phys.\ B}\ \textbf{299}, 719 (1988).

\bibitem{Rybakov2016}
Yu.~P.~Rybakov,
``Topological Solitons in the Skyrme--Faddeev Spinor Model and
Quantum Mechanics,''
\emph{J.\ Phys.:\ Conf.\ Ser.}\ \textbf{731}, 012012 (2016).

\bibitem{Rybakov2024}
Yu.~P.~Rybakov,
``Brioschi 16-Spinors and Photons as Solitons in the Skyrme--Faddeev
Chiral Model,''
\emph{Grav.\ Cosmol.}\ \textbf{30}, 241 (2024).

\bibitem{CODATA2018}
E.~Tiesinga \emph{et al.},
``CODATA Recommended Values of the Fundamental Physical Constants:
2018,''
\emph{Rev.\ Mod.\ Phys.}\ \textbf{93}, 025010 (2021).

\bibitem{Langacker1981}
P.~Langacker,
``Grand Unified Theories and Proton Decay,''
\emph{Phys.\ Rep.}\ \textbf{72}, 185 (1981).

\bibitem{Marciano1981}
W.~J.~Marciano and A.~Sirlin,
``Precise SU(5) Predictions for $\sin^2\theta_W$, $m_W$, and $m_Z$,''
\emph{Phys.\ Rev.\ Lett.}\ \textbf{46}, 163 (1981).

\end{thebibliography}

%% ================================================================
\appendix
\section{Full particle mass spectrum}
\label{app:full_table}
%% ================================================================

Table~\ref{tab:full} lists the complete NWT mass spectrum from
Ref.~\cite{NWT6}, evaluated using Eq.~(\ref{eq:mass}).  Charge
multiplets sharing the same $(p,q,m,n_q)$ are grouped; the
experimental mass shown is the isospin average where applicable.
The 56-particle table from Ref.~\cite{NWT6} includes additional
charge states and achieves 0.40\%~median error; here we show the
distinct mass predictions.

\begin{table}[h!]
\centering
\caption{Full particle mass spectrum from Eq.~(\ref{eq:mass}).}
\label{tab:full}
\smallskip
\small
\begin{tabular}{lcrrrrrrr}
\toprule
Particle & Type & $p$ & $q$ & $m$ & $n_q$
    & $m_\mathrm{exp}$\,(MeV) & $m_\mathrm{pred}$\,(MeV) & Error \\
\midrule
\multicolumn{9}{l}{\textit{Leptons}} \\
$e$            & lepton & 2 & 1 &  3 & 0 &     0.5 &     0.5 & anchor \\
$\mu$          & lepton & 1 & 8 &  9 & 0 &   105.7 &   103.6 & $-2.0\%$ \\
$\tau$         & lepton$^*$ & 3 & 4 & 17 & 3 &  1776.9 &  1789.8 & $+0.7\%$ \\
\midrule
\multicolumn{9}{l}{\textit{Mesons}} \\
$\pi^{\pm}$    & meson  & 3 & 5 &  5 & 2 &   139.6 &   143.3 & $+2.6\%$ \\
$\pi^0$        & meson  & 7 & 3 & 18 & 2 &   135.0 &   134.9 & $-0.1\%$ \\
$K^{\pm}$      & meson  & 2 & 5 &  8 & 2 &   493.7 &   514.8 & $+4.3\%$ \\
$K^0$          & meson  & 7 & 5 & 15 & 2 &   497.6 &   509.1 & $+2.3\%$ \\
$\eta$         & meson  & 6 & 5 & 15 & 2 &   547.9 &   542.7 & $-0.9\%$ \\
$\rho$         & meson  & 5 & 7 &  7 & 2 &   775.3 &   797.2 & $+2.8\%$ \\
$\omega$       & meson  & 4 & 5 & 17 & 2 &   782.7 &   790.9 & $+1.1\%$ \\
$D^{\pm}$      & meson  & 2 & 7 &  5 & 2 &  1869.7 &  1886.2 & $+0.9\%$ \\
$D^0$          & meson  & 3 & 7 &  7 & 2 &  1864.8 &  1844.8 & $-1.1\%$ \\
$J/\psi$       & meson  & 2 & 7 &  7 & 2 &  3096.9 &  3122.9 & $+0.8\%$ \\
$\Upsilon$     & meson  & 4 & 9 &  8 & 2 &  9460.3 &  9434.1 & $-0.3\%$ \\
\midrule
\multicolumn{9}{l}{\textit{Baryons}} \\
$p$            & baryon & 1 & 4 &  5 & 3 &   938.3 &  1032.5 & $+10.0\%$ \\
$n$            & baryon & 1 & 4 &  5 & 3 &   939.6 &  1032.5 & $+9.9\%$ \\
$\Sigma^+$     & baryon & 1 & 4 &  6 & 3 &  1189.4 &  1310.9 & $+10.2\%$ \\
$\Sigma^0$     & baryon & 1 & 4 &  6 & 3 &  1192.6 &  1310.9 & $+9.9\%$ \\
$\Sigma^-$     & baryon & 1 & 4 &  6 & 3 &  1197.4 &  1310.9 & $+9.5\%$ \\
$\Lambda$      & baryon & 3 & 4 & 12 & 3 &  1115.7 &  1123.4 & $+0.7\%$ \\
$\Delta$       & baryon & 5 & 4 & 15 & 3 &  1232.0 &  1222.4 & $-0.8\%$ \\
$\Xi^0$        & baryon & 5 & 4 & 16 & 3 &  1314.9 &  1344.0 & $+2.2\%$ \\
$\Xi^-$        & baryon & 5 & 4 & 16 & 3 &  1321.7 &  1344.0 & $+1.7\%$ \\
$\Sigma^*(1385)$ & baryon & 3 & 4 & 14 & 3 & 1385.0 & 1385.0 & $-0.0\%$ \\
$\Omega^-$     & baryon & 7 & 4 & 19 & 3 &  1672.5 &  1665.6 & $-0.4\%$ \\
\midrule
\multicolumn{9}{l}{\small $^*$See Table~\ref{tab:spectrum} footnote.
    25~distinct predictions; median $|$error$| = 1.1\%$, mean $= 3.0\%$, max $= 10.2\%$.} \\
\multicolumn{9}{l}{\small Full 56-particle table with all charge states:
    median $0.40\%$, from Ref.~\cite{NWT6}.} \\
\bottomrule
\end{tabular}
\end{table}

\end{document}
